\section{Related Work}
\label{sec:related}

\para{Clarification: when to ask.} \citet{zhang2023clarify} give the canonical decomposition ---
when to clarify, what to ask, and how to use the answer --- and propose \textsc{Intent-Sim}, which
scores ambiguity by the entropy of simulated user intents. Their evaluation protocol, AUROC plus
performance under a fixed interaction budget, is the one we adopt in \secref{sec:results} for the
scalar regime. Their absolute numbers are also a caution: AUROCs cluster at 0.50--0.63, and
\textsc{Self-Ask} scores 0.371 on machine translation --- below chance. \citet{zhang2024clamber}
supply a balanced ask/don't-ask benchmark with a taxonomy, which we use as our out-of-source
transfer split precisely because it was never used for tuning. \citet{qian2024in3} annotate the
\emph{importance} of each missing detail, which gives us item-level stakes that nobody announces to
the model. Unlike this line of work, we do not ask whether to clarify; we ask which of four actions
is right, because \emph{don't clarify} is three different things.

\para{Abstention: when to decline.} \citet{wen2024knowlimits} survey abstention along query, model,
and human-value axes and argue it should be treated as a meta-capability rather than a per-task
trick. \textsc{AbstentionBench} \citep{abstentionbench2025} reports that reasoning fine-tuning
\emph{hurts} abstention by roughly 24\%, and ships a human-validated binary abstention judge whose
prompt we extend into a four-way judge (\secref{sec:method}). The limitation we inherit from this
literature is exactly the one it is built on: abstention is binary, so refusing and deferring
score identically.

\para{Agentic help-seeking.} \textsc{HiL-Bench} \citep{trinh2026hilbench} places human-validated
blockers in \textsc{SWE-Bench Pro} and \textsc{BIRD} tasks that surface only through execution, and
finds frontier models score 75--89\% pass@3 with full information but 4--24\% when they must judge
whether to ask --- despite having an \texttt{ask\_human()} tool. Their ASK-F1 metric, the harmonic
mean of question precision and blocker recall, is spam-resistant by construction and we report it
throughout. They also show RLVR on shaped ASK-F1 transfers across domains, which is the strongest
existing evidence for the trainability clause; our \secref{sec:results} training arm supervises the
typed label instead, and finds transfer does \emph{not} follow for free.
\citet{ask2024awn} independently find that agents fabricate missing tool arguments rather than
asking, which is the mechanism behind the over-commitment rates we measure.

\para{The full action space.} Two recent papers propose ontologies close to ours.
\citet{unlu2026ssta} define four support states by minimal repair distance --- complete,
clarifiable, support-blocked, unsupported-now --- and report that scalar confidence prompting
avoids over-commitment but collapses the three-way deferral space to 58.3\% typed accuracy, while
surfacing the categorical ontology reaches 91.7\%. The result is the direct ancestor of our
Experiment~2, but it rests on 32 items, one model, and a dual-persona audit run inside a single
context window, which the author concedes yields upper-bound estimates. We replicate the
\emph{direction} at 15$\times$ the sample size, across five models, with contexts isolated --- and
find the 91.7\% figure was indeed an upper bound (we see 0.650--0.764 typed-deferral accuracy).
\citet{baan2026bag} let a model reason over its own sampled belief state to pick answer, clarify,
or abstain, and name clarify-versus-abstain separation as the open problem; our confusion matrices
quantify it (26\% of \ask items routed to \refuse under typed prompting).
\textsc{Safety Sentry} \citep{safetysentry2026} routes over execute/ask/refuse with a single
decoding-time threshold repositioned per deployment risk --- the cleanest existing mechanism for
the stakes clause, and the motivation for measuring criterion shift separately from discrimination.

\para{Recognition \versus behaviour.} \citet{su2026knowing} show models identify ambiguity when
explicitly asked to judge it, yet default to direct answers in ordinary QA, and that retrieved
context widens the gap. This is the framing of our Experiment~1. Building on it, we do not stop at
showing the gap exists: we decompose recognition into four separately-scored properties and show
that safety and capability recognition are near-ceiling while information-sufficiency recognition
is near chance, which changes what an intervention should target.

\para{Uncertainty estimation.} Semantic entropy over NLI-clustered samples \citep{kuhn2023semantic} is the standard strong
baseline, and conformal methods \citep{yadkori2024conformal} give distribution-free risk control at
a chosen level. \citet{entropyinsufficient2026} report a negative result for entropy-based
selective prediction. \citet{zong2026icalm} announce answer/abstain payoff matrices in the prompt,
which we adapt directly for our announced-stakes manipulation. Our contribution here is negative
and geometric: we show that the problem with a scalar is not its sharpness but its dimensionality
--- three deferral types cannot be recovered from an ordering of one number, and empirically the
models do not even order them consistently.

\para{Evaluation beyond a single turn.} \citet{regretbench2026} score clarification as a
\emph{policy} under hidden intent, decomposing regret against a reference policy into intent
resolution, interaction cost, and ineffective clarification, and \citet{humanevalcomm2024} measure
the cost of a wrong assumption objectively by running the resulting code. Both are the natural
next step past what we measure: \carb scores a single routing decision per request, so it can say
whether asking was right but not whether the conversation that followed was efficient.
